\section{Four New Testament Controls}

\subsection*{The words look toward an approaching death}

The institution narratives do not present the meal as relocation to a future
event. They present Jesus interpreting and giving his approaching death.
Matthew and Mark join the cup to covenant blood poured out ``for many''; Luke
joins the body ``given for you'' and the cup poured out to the command of
memorial (Matt 26:26--29; Mark 14:22--25; Luke 22:14--20). The future Passion
already determines the gift's sacrificial meaning. Yet the narrative continues
from supper to Gethsemane, arrest, trial, and execution. It does not collapse
these events into one historical instant.

\subsection*{Paul holds institution, death, and expectation together}

Paul reports the tradition received from the Lord: the bread and cup are
Christ's Body and new-covenant Blood; the action is done in his memory; and
eating and drinking proclaim the Lord's death until his coming (1 Cor
11:23--26). The celebration therefore has a threefold reach. It obeys the
institution given on the night of betrayal, proclaims the accomplished death,
and awaits the future coming. Paul's grammar preserves the distinction among
institution, accomplished death, and awaited coming. In the rite Christ
sacramentally makes his once-for-all Pasch present and applies its fruits while
the Church still lives between Cross and Parousia.

\subsection*{Hebrews forbids repetition, not participation}

Hebrews insists that Christ entered the sanctuary by his own blood and offered
himself once for all (Heb 7:23--28; 9:11--28; 10:10--14). This excludes a
series of new deaths or supplementary immolations. It does not require the
Cross to be a vanished object available only to recollection. Christ's
priesthood endures, he lives to intercede, and his self-offering has permanent
efficacy. The Catholic account of the Mass must therefore hold together the
historical \emph{once} and the priestly \emph{always}. Literal return to the
Friday of the Passion is not needed to do so.

\subsection*{Biblical memorial is covenantal action}

Catholic tradition reads Christ's memorial command within Israel's Passover
grammar. The Passover memorial does not tell each generation that it has become
the generation physically leaving Egypt. It places each generation within the
covenant identity and saving claim of the Exodus (Exod 12:14, 24--27; 13:3--10;
Deut 5:2--3). In the new covenant the Church teaches that memorial is not bare
recollection: Christ himself and his Pasch are sacramentally made present
(CCC 1363--1364). This exceeds subjective remembering without becoming a
theory of chronological transport.

